\section{软件设计}

软件完整复用学长 \texttt{操作系统和用户程序/} 目录（BootLoader、$\mu$C/OS-II 内核、移植层、Makefile、链接脚本），唯一实质改动在 HAL 文件 \texttt{ucosii/common/openmips.c}：(1) 新增 \texttt{uart\_putc\_raw}；(2) 把 \texttt{TaskStart} 函数体替换为贪吃蛇逻辑。\texttt{main} 函数（\texttt{OSInit $\to$ uart\_init $\to$ gpio\_init $\to$ OSTaskCreate $\to$ OSStart}）保持不变。

\subsection{新增原始字节发送 \texttt{uart\_putc\_raw}}

原 \texttt{uart\_putc} 在发送 \texttt{0x0A}（\texttt{\textbackslash n}）时会自动追加 \texttt{0x0D}（\texttt{\textbackslash r}），这本是为文本换行设计的友好行为，但会\textbf{破坏二进制协议帧的对齐}（位图字节中可能恰好出现 0x0A）。为此新增一个不做任何字符替换的原始发送函数，专用于协议帧：

\begin{lstlisting}[style=c, caption={原始字节发送（openmips.c）}]
void uart_putc_raw(char c) {
    unsigned char lsr;
    WAIT_FOR_THRE;
    REG8(UART_BASE + UART_TH_REG) = c;
    WAIT_FOR_XMITR;
}
\end{lstlisting}

\subsection{贪吃蛇游戏逻辑}

游戏在 16$\times$16 网格上进行。蛇身用文件级静态数组 \texttt{sx[]/sy[]} 存储坐标（下标 0 为蛇头），最大长度 256。游戏状态机有四态：\texttt{READY}（等待开始）、\texttt{RUNNING}（运行）、\texttt{GAMEOVER}（撞墙或撞自身）、\texttt{WIN}（蛇填满全屏）。由于任务栈仅 1\,KB，2\,KB 的蛇身数组必须放在静态全局区而非任务局部变量。

\begin{lstlisting}[style=c, caption={TaskStart 主循环（节选）}]
void TaskStart(void *pdata) {
    OSInitTick();          /* 启动 100Hz 时钟节拍 */
    snake_init();
    while (1) {
        if (state == ST_RUNNING) {
            INT32U b = gpio_in();
            /* 读方向按钮，禁止 180 反向 */
            if ((b & BTN_UP)    && dir != DIR_DOWN)  dir = DIR_UP;
            else if ((b & BTN_DOWN)  && dir != DIR_UP)    dir = DIR_DOWN;
            else if ((b & BTN_LEFT)  && dir != DIR_RIGHT) dir = DIR_LEFT;
            else if ((b & BTN_RIGHT) && dir != DIR_LEFT)  dir = DIR_RIGHT;
            step_snake();     /* 前进 + 碰撞/吃食判定 */
            send_frame();     /* 发状态帧给 PC */
        } else {
            send_frame();
            if (gpio_in() & BTN_CENTER) { snake_init(); state = ST_RUNNING; }
        }
        OSTimeDly(SPEED);     /* SPEED=20 -> 200ms/格 */
    }
}
\end{lstlisting}

\textbf{回环边界}：与经典贪吃蛇"撞墙即死"不同，本实现采用\textbf{回环边界}——蛇从一侧出去后从对侧回来（如左边界出则右边界进），降低难度、提升连贯性：

\begin{lstlisting}[style=c, caption={回环边界处理（step\_snake）}]
if (nx < 0) nx = GW - 1;  else if (nx >= GW) nx = 0;
if (ny < 0) ny = GW - 1;  else if (ny >= GW) ny = 0;
\end{lstlisting}

撞自身（蛇头命中除尾节外的身体）仍判 GAMEOVER；吃到食物则蛇身增长、分数 +1 并用 LCG 伪随机算法在空格上重新投放食物。

\subsection{UART 协议帧}

每次移动打包一个固定 41 字节的二进制帧发送，全量携带蛇身位图，丢帧不影响（下帧覆盖），无需应答。帧格式如表~\ref{tab:frame}。

\begin{table}[H]
    \centering
    \caption{UART 协议帧格式（41 字节，FPGA $\rightarrow$ PC）}
    \label{tab:frame}
    \begin{tabular}{cccp{6.2cm}}
        \toprule
        \textbf{偏移} & \textbf{字段} & \textbf{字节} & \textbf{说明} \\
        \midrule
        0--1 & HEAD & 2 & 帧头 \texttt{0xAA 0x55} \\
        2 & state & 1 & 3=READY / 0=RUNNING / 1=GAMEOVER / 2=WIN \\
        3--4 & score & 2 & 分数，小端 \\
        5 & food & 1 & 食物坐标 \texttt{(fx<<4)|fy} \\
        6 & head & 1 & 蛇头坐标 \texttt{(hx<<4)|hy} \\
        7 & length & 1 & 蛇身长度（含头） \\
        8--39 & bitmap & 32 & 16$\times$16 位图，蛇身格=1，行优先、高位在前 \\
        40 & checksum & 1 & bytes[2..39] 逐字节异或 \\
        \bottomrule
    \end{tabular}
\end{table}

发送时逐字节调用 \texttt{uart\_putc\_raw}（绝不能用 \texttt{uart\_putc}，否则位图中的 0x0A 会引发帧错位）。41 字节/帧在 4800\,bps 下约耗时 86\,ms，200\,ms 的移动节拍足以发完。

\subsection{交叉编译}

因本机与课程均未提供 MIPS 工具链，从 GitHub \texttt{RT-Thread/toolchains-ci} 下载 \texttt{mips-2016.05-7-mips-sde-elf-i686-pc-linux-gnu} 工具链安装到 WSL Ubuntu。该工具链 GCC 版本（5.3）较学长使用的 4.3 新，故不能复用学长预编译的 \texttt{.o}，必须 \texttt{make clean} 后从源码全量重编。编译流程封装在 \texttt{software/build.sh}，最终由 \texttt{BinMerge.exe} 合成 \texttt{OS.bin}（约 32\,KB）。
